\section{Meeting-Success Questionnaire (Davison Instrument)}
\label{annex:davison-questionnaire}

This annex documents the questionnaire used to assess perceived meeting quality in the
user study (Chapter~5).
The instrument is the revalidated meeting-success questionnaire by Davison
\cite{davisonInstrumentMeasuringMeeting1999}.
Item codes in parentheses refer to the original scales:
communication~(C), quality of discussion~(D), status effects~(S),
teamwork~(T), and efficiency~(E).
Unless noted otherwise, items are answered on a five-point scale.

Participants in the user study were German-speaking.
The wording below is therefore the German version administered in the study,
translated from the original English instrument.

The original technology items were written for Group Support Systems.
Additional scaled items on the proactive AI agent, together with two
free-text questions, were therefore appended for this thesis, as discussed
in Chapter~5.
Items A2, A4, and A4b are reverse-coded.

\newcommand{\qbox}{\fbox{\rule{0pt}{1.05ex}\hspace{1.05ex}}}
\newcommand{\qboxes}{\qbox\ \qbox\ \qbox\ \qbox\ \qbox}
\newcommand{\qscale}[2]{%
  {\small #1 \enspace \qboxes{} \enspace #2}%
}
\newcommand{\qitem}[2]{%
  \par\addvspace{0.85em}%
  \noindent #2\hfill~(#1)\\[0.3em]%
}
\newcommand{\qfreelines}{%
  \par\vspace{0.15em}%
  \noindent\rule{\linewidth}{0.4pt}\\[1.15em]%
  \noindent\rule{\linewidth}{0.4pt}\\[1.15em]%
  \noindent\rule{\linewidth}{0.4pt}\\[1.15em]%
  \noindent\rule{\linewidth}{0.4pt}%
}

\begingroup
\small

\subsection*{Eigene Beteiligung}

\noindent
Bezogen auf Ihre eigene Beteiligung am Meeting: Inwieweit stimmen Sie den
folgenden Aussagen zu?

\qitem{C1}{Die Sprache des Meetings hinderte Sie daran, sich zu beteiligen.}
\qscale{Stimme voll zu}{Stimme überhaupt nicht zu}

\qitem{C2}{Es fiel Ihnen schwer, die anderen Gruppenmitglieder zu verstehen,
wenn diese sprachen.}
\qscale{Stimme voll zu}{Stimme überhaupt nicht zu}

\qitem{C3}{Sie hatten Schwierigkeiten, sich auszudrücken.}
\qscale{Stimme voll zu}{Stimme überhaupt nicht zu}

\qitem{C4}{Sie zögerten, eigene Ideen einzubringen.}
\qscale{Stimme voll zu}{Stimme überhaupt nicht zu}

\qitem{S4}{Sie empfanden Druck, entweder einer bestimmten Sichtweise zuzustimmen
oder anderen nicht zu widersprechen.}
\qscale{Stimme voll zu}{Stimme überhaupt nicht zu}

\subsection*{Diskussionsqualität}

\noindent
Bezogen auf alle Meetingteilnehmer:innen insgesamt: Wie bewerten Sie die
Diskussionen im Meeting anhand der folgenden Skalen?

\qitem{D1}{Sinnvoll \enspace \qboxes{} \enspace Sinnlos}

\qitem{D2}{Angemessen \enspace \qboxes{} \enspace Unangemessen}

\qitem{D3}{Offen \enspace \qboxes{} \enspace Geschlossen}

\qitem{D4}{Einfallsreich \enspace \qboxes{} \enspace Einfallslos}

\subsection*{Zusammenarbeit, Effizienz und Statuseffekte}

\noindent
Inwieweit stimmen Sie den folgenden Aussagen zu?

\qitem{T1}{Andere Mitglieder wirkten bereit, Fragen zu beantworten, wenn sie
gestellt wurden.}
\qscale{Stimme voll zu}{Stimme überhaupt nicht zu}

\qitem{T2}{Die Mitglieder arbeiteten als Team zusammen.}
\qscale{Stimme voll zu}{Stimme überhaupt nicht zu}

\qitem{T3}{Die Mitglieder hatten ausreichenden Zugang zu den Informationen, die
  sie benötigten, um sich aktiv zu beteiligen und das Meeting vollständig zu
verstehen.}
\qscale{Stimme voll zu}{Stimme überhaupt nicht zu}

\qitem{E2}{Die im Meeting verbrachte Zeit wurde effizient genutzt.}
\qscale{Stimme voll zu}{Stimme überhaupt nicht zu}

\qitem{E3}{Im Meeting aufgeworfene Themen wurden gründlich diskutiert.}
\qscale{Stimme voll zu}{Stimme überhaupt nicht zu}

\qitem{S1}{Einige Gruppenmitglieder versuchten, andere einzuschüchtern, z.\,B.\
durch lautes Sprechen, aggressive Gesten, Drohungen usw.}
\qscale{Stimme voll zu}{Stimme überhaupt nicht zu}

\qitem{S2}{Einige Gruppenmitglieder versuchten, ihren Einfluss, ihren Status
  oder ihre Macht zu nutzen, um den anderen Gruppenmitgliedern Themen
aufzuzwingen.}
\qscale{Stimme voll zu}{Stimme überhaupt nicht zu}

\qitem{S3}{Sie fühlten sich aufgrund des Verhaltens anderer
  Meetingteilnehmer:innen davon abgehalten, sich an der Diskussion zu
beteiligen.}
\qscale{Stimme voll zu}{Stimme überhaupt nicht zu}

\subsection*{Effizienz und Gesamteindruck}

\qitem{E4}{Wie viel Prozent der Meetingzeit wurden Ihrer Einschätzung nach mit
ernsthafter Diskussion verbracht?}
{\small \underline{\hspace{2.2cm}}\,\%}

\qitem{E1}{In welchem Maße war dieses Meeting ergebnisorientiert?}
\qscale{Stark ergebnisorientiert}{Schwach ergebnisorientiert}

\par\addvspace{0.85em}
\noindent
Wie bewerten Sie Ihre allgemeine Zufriedenheit mit dem Meeting?\\[0.3em]
\qscale{Sehr zufrieden}{Sehr unzufrieden}

\par\addvspace{0.85em}
\noindent
In welchem Maße wurde im Meeting ein Konsens erreicht?\\[0.3em]
\qscale{Stark erreicht}{Schwach erreicht}

\subsection*{Ursprüngliche Technologie-Items (Group Support Systems)}

\par\addvspace{0.85em}
\noindent
Wie wohl fühlen Sie sich im Umgang mit der Technologie?\\[0.3em]
\qscale{Sehr wohl}{Sehr unwohl}

\par\addvspace{0.85em}
\noindent
In welchem Maße hat die Technologie Ihre Beteiligung an diesem Meeting
behindert oder erleichtert?\\[0.3em]
\qscale{Stark behindert}{Stark erleichtert}

\subsection*{Zusätzliche Items zum proaktiven KI-Agenten}

\noindent
The following items were added for this thesis and are not part of the original
Davison instrument.

\qitem{A1}{Die Beiträge des KI-Agenten kamen zu passenden Zeitpunkten.}
\qscale{Stimme voll zu}{Stimme überhaupt nicht zu}

\qitem{A2}{Die Beiträge des KI-Agenten haben den Gesprächsfluss gestört.}
\qscale{Stimme voll zu}{Stimme überhaupt nicht zu}

\qitem{A3}{Der KI-Agent hat bezogen auf den aktuellen Stand der Diskussion
passend reagiert.}
\qscale{Stimme voll zu}{Stimme überhaupt nicht zu}

\qitem{A4a}{Der KI-Agent hat zu häufig eingegriffen.}
\qscale{Stimme voll zu}{Stimme überhaupt nicht zu}

\qitem{A4b}{Der KI-Agent hat zu selten eingegriffen.}
\qscale{Stimme voll zu}{Stimme überhaupt nicht zu}

\qitem{A5}{Es war in Ordnung, dass der KI-Agent von sich aus gesprochen hat,
ohne gefragt zu werden.}
\qscale{Stimme voll zu}{Stimme überhaupt nicht zu}

\qitem{A6}{Der KI-Agent hat dazu beigetragen, dass Informationen zur Sprache
kamen, die sonst unerwähnt geblieben wären.}
\qscale{Stimme voll zu}{Stimme überhaupt nicht zu}

\qitem{A7}{Die Hinweise des KI-Agenten waren so formuliert, dass niemand
bloßgestellt wurde.}
\qscale{Stimme voll zu}{Stimme überhaupt nicht zu}

\subsection*{Ergänzende Freitextfragen}

\noindent
The following open questions were appended in order to
collect qualitative accounts of how participants experienced the agent's
contributions, as discussed in Chapter~5.

\qitem{F1}{Beschreiben Sie eine Situation, in der der KI-Agent eingegriffen hat.
  Wie haben Sie diesen Beitrag erlebt, und welchen Einfluss hatte er auf die
Diskussion?}
\qfreelines

\qitem{F2}{Was war an den Beiträgen des KI-Agenten hilfreich und was hat eher
gestört oder gefehlt?}
\qfreelines

\endgroup
